% Terjemahan Bahasa Indonesia dari prelude.tex baris 154-183.
% Sumber: Methods of Algebra, Volume 2, commit
% 9a5803ff77dd3257484cb177f851a73770a59dd3 (CC BY 4.0).
% stable-unit-id: o014.aljabr2.prelude.structure-overview

% segment-id: o014.li2.prelude.structure-overview.h001
\section*{Ikhtisar Struktur}
\addcontentsline{toc}{section}{Ikhtisar Struktur}
\hypertarget{o014.aljabr2.prelude.structure-overview}{ }

% segment-id: o014.li2.prelude.structure-overview.p001
Pembaca disarankan memulai dari bagian Konvensi pada akhir pendahuluan.
Tujuannya ialah mengenali terlebih dahulu simbol dan konvensi yang
digunakan di seluruh buku. Selain itu, pembaca pemula tidak dianjurkan
mengikuti buku ini secara ketat dari awal hingga akhir, kecuali jika sudah
cukup siap menerima metode abstrak. Awal setiap bab menyediakan ikhtisar
dan petunjuk bacaan yang terperinci.

% segment-id: o014.li2.prelude.structure-overview.p002
Bagian pertama ialah \emph{Bagian Dalam}, yang meletakkan landasan.

\begin{description}
	% segment-id: o014.li2.prelude.structure-overview.i001
	\item[Bab 1: Pelengkap Teori Kategori] Mengulas dan melengkapi pengetahuan
	teori kategori yang diperlukan, dengan pokok utama berupa morfisme ketat
	(\emph{strict morphism}), kernel dan kokernel dalam kategori aditif,
	ekstensi Kan (\emph{Kan extension}), lokalisasi Gabriel--Zisman, dan topik-topik lain.
	% segment-id: o014.li2.prelude.structure-overview.i002
	\item[Bab 2: Kategori Abelian] Membangun teori dasar kategori abelian,
	dilengkapi dengan materi mengenai kategori Grothendieck dan topik lain.
	% segment-id: o014.li2.prelude.structure-overview.i003
	\item[Bab 3: Kompleks] Membahas konsep-konsep dasar berupa kompleks,
	bikompleks, homotopi, kerucut pemetaan, dan sebagainya pada kategori
	aditif, serta kohomologi dalam kasus kategori abelian. Paruh kedua bab ini
	membahas funktor turunan dari sudut pandang klasik beserta kasus-kasus
	khusus seperti $\Ext$, $\Tor$, dan $\lim^1$. Untuk menangani kompleks tak
	terbatas, bagian akhirnya juga mengkaji resolusi $K$-injektif dan
	$K$-projektif.
	% segment-id: o014.li2.prelude.structure-overview.i004
	\item[Bab 4: Kategori Bertriangulasi dan Kategori Turunan] Memperkenalkan
	konsep umum kategori bertriangulasi dan funktor bertriangulasi, lalu
	berfokus pada kategori turunan dari kategori abelian serta funktor
	turunan antara kategori-kategori turunan. Dari sudut pandang klasik,
	$\Ext$ dan $\Tor$ kemudian diangkat menjadi $\RHom$ dan $\otimesL$.
	% segment-id: o014.li2.prelude.structure-overview.i005
	\item[Bab 5: Barisan Spektral] Berangkat dari definisi umum barisan
	spektral dan pasangan eksak, lalu secara bertahap mengkhususkan
	pembahasan pada barisan spektral kompleks terfilter dan bikompleks,
	beserta penerapannya dalam berbagai konteks aljabar. Pasangan eksak tidak mutlak
	diperlukan untuk barisan spektral kompleks terfilter.
\end{description}

% segment-id: o014.li2.prelude.structure-overview.p003
Berikutnya ialah \emph{Bagian Luar}, yang memuat penerapan atau, dengan kata
lain, pengembangan lebih lanjut.
\begin{description}
	% segment-id: o014.li2.prelude.structure-overview.i006
	\item[Bab 6: Homologi dan Kohomologi Grup] Mengkaji homologi dan
	kohomologi grup beserta berbagai perhitungan atau contoh terkait. Bab ini
	juga membahas kohomologi Tate, kohomologi grup profinit, dan kohomologi
	nonabelian (\emph{nonabelian cohomology}), yang masing-masing memegang peranan penting dalam
	penerapan.
	% segment-id: o014.li2.prelude.structure-overview.i007
	\item[Bab 7: Teori Monad] Berangkat dari konsep umum ``aljabar'' dalam
	kategori monoidal dan mula-mula membahas struktur diferensial bergradasi
	yang berkaitan dengan kompleks. Selanjutnya, setelah menguraikan contoh
	seperti koaljabar, bialjabar, dan aljabar Hopf, pembahasan beralih ke
	teorema monadisitas Beck. Bagian kedua bab mengkaji teori modul,
	memperkenalkan teori Morita, lalu menggunakan teorema monadisitas untuk
	menangani penurunan datar (\emph{flat descent}) bagi
	modul dan penurunan Galois (\emph{Galois descent}).
	% segment-id: o014.li2.prelude.structure-overview.i008
	\item[Bab 8: Metode Simplisial] Mencakup teori umum objek simplisial dan
	himpunan simplisial, korespondensi Dold--Kan, serta konstruksi bar yang
	berlandaskan bahasa monad. Subbab-subbab terakhir membahas teorema
	Eilenberg--Zilber yang diperumum dan penafsiran topologis atas kerucut
	pemetaan.
	% segment-id: o014.li2.prelude.structure-overview.i009
	\item[Bab 9: Dualitas] Berangkat dari dualitas dalam kategori monoidal,
	lalu, setelah memberikan contoh-contoh dasar, mengkaji teorema
	rekonstruksi dan kategori Tannakian. Isi bab ini tidak berhubungan
	langsung dengan kompleks.
\end{description}

% segment-id: o014.li2.prelude.structure-overview.p004
Menurut rencana semula, \emph{Bagian Luar} seharusnya mencakup materi dasar
tentang teori representasi, yang diperkirakan akan berperan secara alami dalam
keseluruhan susunan buku. Pada tahap akhir, rencana tersebut dibatalkan karena
buku ini sudah terlalu panjang dan telah tersedia beberapa buku ajar yang
mapan.

% segment-id: o014.li2.prelude.structure-overview.p005
Isi lampiran dirinci sebagai berikut.
\begin{description}
	% segment-id: o014.li2.prelude.structure-overview.i010
	\item[Lampiran A: Materi Lanjutan tentang Kategori Abelian] Memuat bahan
	yang berkaitan secara langsung atau tidak langsung dengan kategori
	abelian, baik yang dirujuk oleh teks utama maupun yang merupakan materi
	sampingan.
	% segment-id: o014.li2.prelude.structure-overview.i011
	\item[Lampiran B: Pengantar Objek-ind dan Objek-pro] Berangkat dari teori
	grup profinit, lalu secara alami memperkenalkan konsep objek-ind dan
	objek-pro serta mengkaji konstruksi Ind pada kategori dan funktor.
	Setelah itu, dalam kasus khusus kategori abelian, lampiran ini membuktikan
	teorema pembenaman Freyd--Mitchell sebagai contoh penerapan.
\end{description}

% segment-id: o014.li2.prelude.structure-overview.p006
Terakhir, laman pribadi penulis menyediakan informasi tambahan terkini,
termasuk daftar ralat, baik untuk buku ini maupun karya lainnya. Pembaca
dipersilakan memantau laman tersebut. Jika menemukan kesalahan dalam buku ini,
pembaca dimohon mengirimkan koreksi kepada penulis.
